\section*{Paper \ref{paper:slr}}
\textbf{\citetitle{miticSystematicLiteratureReview2026}}

\subsubsection*{Authors}
\citefullauthors{miticSystematicLiteratureReview2026}

\subsubsection*{Status}
Accepted to \citepublication{miticSystematicLiteratureReview2026}, May 2026

\subsubsection*{Research methods}
\gls{slr} following the \gls{prisma} standard to gather all currently available methods for determining spatial features of stereophonic musical audio.

\subsubsection*{Results and conclusions}
Addressing \rqb{}, the review revealed a comprehensive set of techniques that can be used to make spatial descriptions of stereo music, encompassing low-level spatial cues, high-dimensional descriptors, and predictions of important perceptual values. Despite the high number of studies that were assessed, there were relatively few that were \enquote{very relevant} to recorded music specifically. The \gls{slr} also revealed that while \gls{ml}/\gls{dl} have been used on spatial features in the past, both \textit{unsupervised} \gls{ml} and very large music corpora remain fairly unexplored. It also details a framework for how to create a \gls{sip} algorithm that is useful for any given dataset and use case. It acts as a one-stop-shop to springboard any future research into stereo field.

\subsubsection*{My contribution}
Main author, responsible for all data collection and processing, and writing.